\chapter{Safety-Critical Systems, RE Tooling \& Integration}
\section{Safety Risk Analysis}
Each hazard records a cause, consequence, severity S, and likelihood L. This report uses $R=S\times L$ to prioritize review within the sample dataset. S has four levels and L has three; the workbook does not quantify the meaning of each level, so scores support only internal case-study comparisons.
\par\needspace{9\baselineskip}\begingroup\footnotesize\setlength{\tabcolsep}{4pt}
\begin{longtable}{>{\raggedright\arraybackslash}p{\dimexpr 0.10000\linewidth-2\tabcolsep\relax}>{\raggedright\arraybackslash}p{\dimexpr 0.30000\linewidth-2\tabcolsep\relax}>{\raggedright\arraybackslash}p{\dimexpr 0.30000\linewidth-2\tabcolsep\relax}>{\raggedright\arraybackslash}p{\dimexpr 0.30000\linewidth-2\tabcolsep\relax}}
\caption{Risk matrix S x L}\label{tab:risk-grid}\\
\toprule
\textbf{S / L} & \textbf{1} & \textbf{2} & \textbf{3} \\ \midrule
\endfirsthead
\multicolumn{4}{l}{\small\itshape Risk matrix S x L (continued)}\\
\toprule
\textbf{S / L} & \textbf{1} & \textbf{2} & \textbf{3} \\ \midrule
\endhead
\midrule
\multicolumn{4}{r}{\scriptsize Continued on the next page}\\
\endfoot
\bottomrule\endlastfoot
4 & HZ-04, HZ-07 & HZ-01 & --- \\
3 & --- & HZ-05, HZ-10 & HZ-06, HZ-11 \\
2 & --- & HZ-09 & HZ-02, HZ-03, HZ-08, HZ-12 \\
1 & --- & --- & --- \\
\end{longtable}\endgroup

\par\begingroup\footnotesize\setlength{\tabcolsep}{4pt}
\begin{longtable}{>{\raggedright\arraybackslash}p{\dimexpr 0.22000\linewidth-2\tabcolsep\relax}>{\raggedright\arraybackslash}p{\dimexpr 0.10000\linewidth-2\tabcolsep\relax}>{\raggedright\arraybackslash}p{\dimexpr 0.08000\linewidth-2\tabcolsep\relax}>{\raggedright\arraybackslash}p{\dimexpr 0.38000\linewidth-2\tabcolsep\relax}>{\raggedright\arraybackslash}p{\dimexpr 0.22000\linewidth-2\tabcolsep\relax}}
\caption{Hazard register and mitigation scope (ordered by R)}\label{tab:hazards}\\
\toprule
\textbf{ID / hazard} & \textbf{S / L / R} & \textbf{SIL} & \textbf{Mitigation and limitations} & \textbf{Traceability} \\ \midrule
\endfirsthead
\multicolumn{5}{l}{\small\itshape Hazard register and mitigation scope (ordered by R) (continued)}\\
\toprule
\textbf{ID / hazard} & \textbf{S / L / R} & \textbf{SIL} & \textbf{Mitigation and limitations} & \textbf{Traceability} \\ \midrule
\endhead
\midrule
\multicolumn{5}{r}{\scriptsize Continued on the next page}\\
\endfoot
\bottomrule\endlastfoot
HZ-06 Plate misrecognition & 3 / 3 / 9 & SIL-1 & 80\% threshold; manual fallback; confidence is simulated. & REQ-11 \\
HZ-11 Network loss stops gates & 3 / 3 / 9 & SIL-1 & Local entry/exit, marked logs; actual network synchronization is absent. & REQ-29; REQ-30; NFR-02 \\
HZ-01 Owner data exposure & 4 / 2 / 8 & SIL-2 & Encrypt backups and restrict access; no encryption evidence. & NFR-04 \\
HZ-02 Incorrect fees & 2 / 3 / 6 & QM & Test time/time-zone boundaries; overrides need manager approval. & REQ-03; REQ-22 \\
HZ-03 Duplicate slot assignment & 2 / 3 / 6 & QM & Reject repeated Occupied transitions; concurrent races are unverified. & REQ-01 \\
HZ-05 Unauthorized zone changes & 3 / 2 / 6 & SIL-1 & Authenticate, lock after three failures, and restrict to admin. & REQ-18; REQ-09 \\
HZ-08 Missing/incorrect sensor data & 2 / 3 / 6 & QM & Reconcile with tickets, alert, and block faulty slots. & REQ-13 \\
HZ-10 Exit with another ticket & 3 / 2 / 6 & SIL-1 & Compare exit and ticket plates; approval and logging on mismatch. & REQ-26 \\
HZ-12 Override abuse & 2 / 3 / 6 & QM & Separate request/approval roles; enforce distinct user identities in future. & REQ-22 \\
HZ-04 Barrier strikes a vehicle & 4 / 1 / 4 & SIL-2 & The gate refuses to close on a vehicle, a sensor fault, or power loss (REQ-31); no physical interlock is tested. & NFR-02; REQ-31 \\
HZ-07 Evacuation barriers fail to open & 4 / 1 / 4 & SIL-2 & Fail-open and payment bypass; physical and power-loss tests are needed. & REQ-15; NFR-06 \\
HZ-09 Excessive dynamic fees & 2 / 2 / 4 & QM & Fixed multipliers 1.5 x 1.25; show breakdown; clarify the cap. & REQ-16; REQ-17; REQ-28 \\
\end{longtable}\endgroup

SIL-1, SIL-2, and QM labels are preserved for source traceability. The \texttt{sil\_rule} function derives a label from severity alone. This is an educational convention, and there is no analysis or evidence supporting a claim of certified safety integrity. PII exposure, fraud, and revenue discrepancies also appear in the hazard register despite differing from physical injury risks.

The simulated policy keeps a barrier open on sensor failure, vehicle presence, or power loss. Separately, evacuation bypasses payment and blocks new tickets. REQ-31 makes this explicit: \texttt{close\_gate} calls \texttt{close\_barrier} and refuses to close on a vehicle, a sensor fault, or power loss (TC-67). The logic is not connected to hardware, so no physical interlock is demonstrated; HZ-04 is therefore covered by a requirement and a software test only. HZ-01 encryption has no verification demonstration. Other hazards have logical demonstrations, but HZ-03 has not been tested under concurrent allocation.

REQ-22 separates roles: an attendant requests an override, and a manager or administrator approves it. The code checks roles, not identities, so one person holding both roles could request and approve an override. This limitation matters if the mechanism is used operationally.

\section{Integrated Requirements Management Tool}
\begin{table}[H]\centering\small\caption{Modules and responsibilities}
\begin{tabularx}{\linewidth}{L{28mm}X}\toprule
File & Responsibility\\\midrule
\texttt{parking.py} & Entry point, menus, interactive mode, and Markdown appendix export.\\
\texttt{rm.py} & Excel loading, quality checks, traceability, dependencies, CR impact, risk matrices, and Mermaid export.\\
\texttt{sim.py} & Lot hierarchy, tickets, states, pricing, authorization, sensors, evacuation, and persistence.\\
\texttt{checks.py} & 22 use-case scenarios, 67 tests, NFR measurements, and hazard-mitigation demonstrations.\\
\texttt{live.py} & Seven end-to-end gate scenarios using interactive input or default answers.\\\bottomrule
\end{tabularx}\end{table}
The workflow starts with Excel, applies data checks and traceability, and executes scenarios on the simulation to collect results. The tool exports Markdown and Mermaid. This report turns the dataset into LaTeX tables and figures with execution evidence. Jira and PlantUML/Mermaid are listed as reference tools in the workbook; no actual Jira connection is evidenced.

\begin{Verbatim}[fontsize=\small,frame=single]
python parking.py
python parking.py all
\end{Verbatim}
The first command opens the menu. The \texttt{all} command runs feature groups, writes test results to the workbook, and exports \texttt{out/report\_appendix.md}. Use a copy when demonstrating it if the baseline must remain unchanged. Interactive mode saves \texttt{state.json}; dynamically added zones are not fully reconstructed after restart.

\subsection*{Sample Tool Output}
The quality check (menu item 2) produced the following findings on the report dataset. The traceability and change-impact outputs are reproduced in Tables~\ref{tab:trace} and~\ref{tab:changes}.
\VerbatimInput[fontsize=\scriptsize,frame=single]{evidence/quality.txt}

\section{Simulation \& End-to-end Testing}
\par\needspace{12\baselineskip}\begingroup\footnotesize\setlength{\tabcolsep}{4pt}
\begin{longtable}{>{\raggedright\arraybackslash}p{\dimexpr 0.32000\linewidth-2\tabcolsep\relax}>{\raggedright\arraybackslash}p{\dimexpr 0.68000\linewidth-2\tabcolsep\relax}}
\caption{Seven recorded live scenarios}\\
\toprule
\textbf{Scenario} & \textbf{Expected observations} \\ \midrule
\endfirsthead
\multicolumn{2}{l}{\small\itshape Seven recorded live scenarios (continued)}\\
\toprule
\textbf{Scenario} & \textbf{Expected observations} \\ \midrule
\endhead
\midrule
\multicolumn{2}{r}{\scriptsize Continued on the next page}\\
\endfoot
\bottomrule\endlastfoot
1. Entry / ALPR & Plate, confidence, slot, ticket, and open/close gate states. \\
2. Exit / payment & A fee breakdown crossing peak hours; total and simulated receipt. \\
3. Reservation to exit & Held slot, arrival within hold, sensor scan, and fee. \\
4. Evacuation drill & Open gates; refused entry; zero-fee exit; administrator ends evacuation. \\
5. Ticket swap & Mismatch held; staff confirmation audited. \\
6. Network outage & Local entry, refused online payment, successful cash exit, and marked-log recovery. \\
7. EV / monthly pass & Energy plus parking; monthly pass waives parking. \\
\end{longtable}\endgroup

The simulated day begins with two cars and a motorcycle and includes low-confidence recognition, a reservation no-show, zone locking, a sensor mismatch, and lost-ticket handling. Logs are produced by the existing code. Labels such as ALPR CAM, PAYMENT, and CHARGER are console labels, not evidence of connections to corresponding equipment.
\begin{Verbatim}[fontsize=\scriptsize,frame=single]
10:00  29A-111 identified by camera, enters -> A1
10:10  29B-222 camera confidence 60% -> manual entry -> A2
10:15  M-1 enters -> B1
10:20  VIP-1 reserves for 11:00 -> C1
10:30  admin locks zone B -> 1 slot(s) blocked
10:40  M-2 enters -> rejected (zone B locked)
10:50  sensor scan -> ['missing_vehicle']
11:20  VIP-1 no-show -> released ['C1']
12:00  29A-111 exits, fee 20000 VND
12:30  29B-222 lost ticket, fee 290000 VND
13:00  M-1 exits, fee 9000 VND

LOT Smart Parking  (4/6 free, revenue 319000)
  FLOOR F1  (2/4 free, revenue 319000)
    ZONE A  (2/2 free, revenue 310000)
      [A1] Free 10m
      [A2] Free 13m
    ZONE B  (0/2 free, revenue 9000)
      [B1] Blocked 25m
      [B2] Blocked 28m
  FLOOR F2  (2/2 free, revenue 0)
    ZONE C  (2/2 free, revenue 0)
      [C1] Free 50m
      [C2] Free 53m

Audit log: 16 records (alert_missing_vehicle, anpr_low_confidence, anpr_read, entry, entry_rejected,
  exit, lock, lost_ticket, manual_entry, payment, reserve, reserve_expired).
\end{Verbatim}

The seven live scenarios add entry/exit, reservation-to-exit journeys, evacuation, ticket swapping, network loss, and EV/monthly-pass operation. Complete output is stored in \texttt{evidence/live\_runs.txt}. Default responses provide repeatable execution and do not introduce new survey data.

\section{Chapter Conclusion}
The tool connects requirement data, models, and experiments in an inspectable workflow. Safety analysis illustrates evacuation priority while identifying the need to validate hardware, interlocks, security, and concurrent data handling before deployment.

The workflow runs from the workbook through quality checks, traceability, change impact, the 22 use-case scenarios, 67 tests, and seven live gate scenarios, and it writes the test results back to the workbook so that the evidence can be inspected. Of the 12 hazards, HZ-04 gained a requirement (REQ-31) and a software test (TC-67) during this work, but no physical interlock is demonstrated. HZ-01 has no encryption demonstration and HZ-03 has no concurrent-allocation test. Override approval checks roles, not identities. These limits define the work that remains before any real deployment.
